\chapter{Building on a public surface}
\label{ch:surface}

\section*{What this chapter is for}

There is one decision that changes an artefact's value more than anything else
in it, and it is made before any code is written: what the artefact is built
\emph{against}.

\section{The move}

Most portfolio work is built against something invented. A fictional domain, a
sample dataset, an imagined API. It demonstrates the technique and nothing
else, and the reader has to take on trust that it would survive contact with
something real.

The alternative is to build against a surface that is already published and
that you do not control: a documented API, an open protocol, a public data
feed, an open-source project's extension points. You become an external
consumer of somebody's published work, and then you publish what happened.

\begin{kitnote}
The surface has to be genuinely public and genuinely outside your control. That
is what supplies the property: a published interface has constraints you did
not choose, documentation that is occasionally wrong, errors that are not the
ones you would have designed, and a shape you have to accommodate rather than
negotiate. Every one of those is a thing an invented surface cannot give you.
\end{kitnote}

\section{What it buys that an invented surface cannot}

\textbf{The findings are real.} You will discover that the documentation and
the behaviour disagree somewhere, that an error case is undocumented, or that
the obvious mapping loses information. Those discoveries are the content.
Nobody can produce them against a surface they designed.

\textbf{It is checkable.} A reader can go and look at the same published
surface. That removes the residual doubt that the interesting part was arranged.

\textbf{It demonstrates the actual job.} Chapter~\ref{ch:three-jobs} put
integration, evaluation and operation at the centre of this role. Building
against something you did not design is that job in miniature.

\textbf{It gives the conversation somewhere to go.} ``What surprised you?'' is
answerable with something specific, and specificity is most of what separates a
good technical conversation from a recital.

\section{Choosing one}

\begin{kitmethod}
\textbf{1. Prefer a surface in the segment you chose} in Chapter~\ref{ch:segment}.
An integration against a developer tool reads differently in an infrastructure
loop than in a product one.

\textbf{2. Prefer documented over popular.} You want a specification to read and
compare against behaviour. Popularity gives you tutorials, which is the
opposite of what you need \dash{} a surface with thirty existing write-ups has
had its interesting findings taken.

\textbf{3. Prefer one with an unglamorous edge.} Pagination, partial failure,
rate limiting, an inconsistent identifier scheme. The finding lives at the edge,
not in the happy path.

\textbf{4. Check you can access it without a commercial relationship.} If a
credential requires a contract, you will end up with the third register class
from Chapter~\ref{ch:ledger} \dash{} an integration that has never run \dash{}
and you should know that going in rather than discovering it in week two.

\textbf{5. Write the thesis before the code.} ``What breaks when you put a model
behind an API that was designed for a person'' is a thesis. ``An integration
with a public API'' is not.
\end{kitmethod}

\section{If you target a company's own surface}

Where the company you are applying to publishes something \dash{} an API, a
protocol implementation, an open-source library \dash{} building against it
before applying is the strongest version of this move, for a reason that has
nothing to do with flattery: you arrive as a user of their work with findings,
rather than as an applicant with a portfolio.

Three cautions, in order of how expensive they are to get wrong.

\textbf{Findings are findings, not criticism.} ``This behaviour is not in the
documentation'' is a gift. ``Your API is badly designed'' is an opinion about
work whose constraints you do not know, delivered to the people who made it.
\judgement{The register from Chapter~\ref{ch:ledger} is what keeps you on the
right side of this: a finding stated with what you did not check is a finding;
the same sentence without it is a verdict.}

\textbf{The artefact must stand up without the company.} If it only makes sense
as an application, it is an application. Write it so it would be worth
publishing if you never applied, and it will be worth more when you do.

\textbf{Publish before you apply, not after.} The sequence is the substance of
the claim. Something published afterwards is a response to a process;
something published before is what you were doing anyway.

\begin{kittrap}
Do not build against a surface you cannot actually reach, and then write it up
as though you had. An integration written against a published specification and
never executed is a legitimate artefact \emph{if the register says so} and a
serious problem if it does not \dash{} because the first specific question will
be about behaviour you have never observed, and there is no recovery from
having implied otherwise.
\end{kittrap}

\begin{drill}
List three public surfaces you could build against, in your chosen segment.
For each, write the one-sentence thesis and the one edge case you expect to be
interesting. Then spend thirty minutes actually calling one of them. The list
usually reorders itself completely once one of the three turns out to need a
commercial agreement and another turns out to be better documented than it
looked.
\end{drill}

\section*{What this chapter does not claim}

It does not claim this produces interviews. It claims the artefact is worth more
per hour spent than the same work against an invented surface, because the
findings are real and checkable.

It makes no claim about how any particular company responds to somebody building
on their published work. That depends entirely on the company and on the tone,
which is why two of the three cautions above are about tone.
